\section{Evaluation}

The evaluation aims to measure both the correctness and
computational efficiency of the proposed neuro-symbolic decoding
architecture. In particular, the experiments evaluate how effectively the
Syntax Oracle reduces structural and semantic errors during generation while
minimizing runtime interventions.

\subsection{Target Domain}
All experiments will focus on structured JSON generation governed by
a specification language with CSI annotations to form our abstract machine.
These annotations express cross-field constraints such as arithmetic relationships (e.g.,
\texttt{max\_items} $\geq$ \texttt{min\_items}), conditional dependencies,
and referential integrity between identifiers.

Evaluation datasets will consist of natural language prompts paired with
target JSON outputs, drawing from benchmarks such as BIRD~\cite{li2023can} for relational 
complexity and specialized schema-following tasks. Samples will be drawn from both Oracle-generated
synthetic schemas and instruction-conditioned tasks to ensure coverage of
diverse structural configurations.

\subsection{Performance Metrics}
System performance will be evaluated using several complementary metrics.

\begin{enumerate}
\item Oracle Intervention Rate (OIR).

The primary metric introduced in this work measures how frequently the
Syntax Oracle must intervene during generation through masking, semantic
validation failure, or rollback. Lower intervention rates indicate that the
model has internalized more of the schema constraints.

To normalize for task difficulty, OIR will be entropy-scaled using the model's 
token distribution $H(p_t)$, a technique adapted from predictive coding 
metrics:

\begin{equation}
\text{OIR}*{norm} =
\frac{\text{Oracle interventions}}
{\sum*{t=1}^{T} H(p_t)}
\end{equation}

\item Semantic Validity Rate.

The proportion of generated outputs that satisfy all schema constraints when
validated by the Oracle, similar to the "Pass@k" metric used in code generation 
benchmarks~\cite{chen2021evaluating}.

\item Inference Efficiency.

Computational overhead will be measured using standard LLM performance
benchmarks, including decoding latency (ms/token) and generation throughput 
(tokens/second), utilizing the benchmarking methodology defined 
in~\cite{miao2024specinfer}.

\end{enumerate}

\subsection{Constraint Complexity Scaling}
To evaluate scalability, experiments will vary schema complexity by
increasing the number of fields, semantic constraints, and dependency depth.
We will specifically analyze the "Semantic Wall" phenomenon 
identified in~\cite{valmeekam2023planning}, quantifying how 
constraint density affects both neural prediction behavior and 
verification workload.

\subsection{Baseline Comparisons}
The proposed system will be compared against several baseline approaches:

\begin{itemize}
\item unconstrained autoregressive generation
\item prompt-based schema enforcement
\item syntax-only FSM masking
\item LoRA fine-tuning without runtime verification
\end{itemize}

These baselines represent the dominant strategies for structured generation.
Additional ablation experiments will isolate the contribution of syntactic
masking, semantic rollback, and Oracle-guided training within the proposed
framework.